\documentclass[11pt]{article}
\usepackage[margin=1in]{geometry}
\usepackage{booktabs}
\usepackage{longtable}
\usepackage{hyperref}
\usepackage{xcolor}
\usepackage{listings}
\usepackage{parskip}
\usepackage{enumitem}

\title{Replication Report --- BVBRC-68: pPA1011 \textit{bla}\textsubscript{KPC-2} Plasmid \\
       \large Hu et al. 2019, \textit{Infect Drug Resist} 12:1285--1288}
\author{Independent replication (free/public data only)}
\date{}

\begin{document}
\maketitle

\begin{center}
\fbox{\parbox{0.9\textwidth}{\centering
\textbf{Verdict:} \textbf{\textcolor{blue}{REPLICATED}} (structural + molecular; on the sequence claims that are checkable from public data).\\[4pt]
\textbf{One-line:} pPA1011 (GenBank MH734334.1) is 62{,}793~bp with GC 58.78\% (paper: 62{,}793~bp, 58.8\%), and its \textit{bla}\textsubscript{KPC-2} gene translates to 293~aa with 100.00\% identity to the canonical KPC-2 protein.
}}
\end{center}

\section{Target paper}
Hu YY, Wang Q, Sun QL, Chen GX, Zhang R. \emph{A novel plasmid carrying carbapenem-resistant gene bla\textsubscript{KPC-2} in Pseudomonas aeruginosa.} Infect Drug Resist. 2019 May 14;12:1285--1288. \texttt{doi:10.2147/IDR.S196390}. PMID: 31190916; PMCID: PMC6526029.

\section{Paper summary}
Hu et al.\ (2019, Zhejiang University) report a carbapenem-resistant \textit{P.~aeruginosa} strain \textbf{PA1011} (MLST \textbf{ST463}) isolated from a surgical ICU patient. PCR confirmed \textit{bla}\textsubscript{KPC-2} carriage. The plasmid, isolated and sequenced (Illumina NextSeq~500 + PacBio~RSII), was named \textbf{pPA1011}. Reported features:
\begin{itemize}[nosep]
  \item Length: \textbf{62{,}793~bp}
  \item GC content: \textbf{58.8\%}
  \item Novel plasmid backbone
  \item Novel \textit{bla}\textsubscript{KPC-2} genetic environment: \textbf{$\Delta$IS6--Tn3--ISKpn8--\textit{bla}\textsubscript{KPC-2}--ISKpn6--IS26}
  \item Deposited at NCBI as \textbf{MH734334.1}.
\end{itemize}

\section{Claims table}
\begin{longtable}{@{}p{0.05\textwidth}p{0.42\textwidth}p{0.20\textwidth}p{0.27\textwidth}@{}}
\toprule
ID & Claim & Public-testable? & Test used \\ \midrule
C1 & pPA1011 is 62{,}793~bp & Yes & Length of MH734334.1 FASTA \\
C2 & GC content 58.8\% & Yes & Recompute GC from FASTA \\
C3 & Carries \textit{bla}\textsubscript{KPC-2} (identical to canonical) & Yes & Translate excised 882-bp ORF; align to canonical KPC-2 (293~aa) \\
C4 & Genetic environment: Tn3 $\to$ ISKpn $\to$ \textit{bla}\textsubscript{KPC-2} $\to$ IS elements & Partially & Parse GenBank features $\pm 6$~kb around \textit{bla}\textsubscript{KPC-2} \\
C5 & Novel backbone (not previously described) & Only in ``novel configuration'' sense & BLAST vs closest published KPC plasmid p14057 (KY296095.1) \\
C6 & PA1011 is ST463 & Provenance-only (no WGS on SRA) & Cross-check MH734334 record note \\
\bottomrule
\end{longtable}

\section{Method}
All analyses were performed inside the replication directory \\
\path{~/Dropbox/REPLICATE-PROJECT/BVBRC-68-pseudomonas-blakpc2-plasmid/}. Free tools only; no paid endpoints used.

\subsection{Data sources}
\begin{enumerate}[nosep]
  \item \textbf{MH734334.1} --- pPA1011 complete plasmid (NCBI GenBank). Depositor: Hu YY et al., Second Affiliated Hospital of Zhejiang University; LOCUS confirms 62{,}793~bp, circular, plasmid = \texttt{p1011-KPC2}, geo = China, note = \texttt{genotype: ST463}.
  \item \textbf{KY296095.1} --- p14057, a Chinese \textit{P.~aeruginosa} \textit{bla}\textsubscript{KPC-2} plasmid used as the closest published KPC-plasmid comparator (51{,}663~bp).
  \item Canonical KPC-2 protein (293~aa) --- reference identical to the beta-lactamase KPC ProteinID AZZ88873.1 embedded in MH734334.1.
\end{enumerate}

\subsection{Analyses}
\begin{itemize}[nosep]
  \item \textbf{Length \& GC (C1, C2):} Python parses the MH734334.1 FASTA; counts A/C/G/T; GC\% $= (G+C)/\text{total}\times 100$.
  \item \textbf{KPC-2 identity (C3):} 6-frame translation of the 882-bp region excised at the \textit{bla}\textsubscript{KPC-2} CDS coordinates; longest M-to-* ORF compared position-by-position to canonical KPC-2.
  \item \textbf{Genetic environment (C4):} Parse GenBank feature table; list all CDS / mobile-element / repeat features within $\pm 6$~kb of \textit{bla}\textsubscript{KPC-2}.
  \item \textbf{Novelty vs prior KPC plasmid (C5):} \texttt{blastn -task megablast} of pPA1011 vs p14057, 19~HSPs; aggregate query-side coverage (union of covered positions on pPA1011) and length-weighted mean \% identity.
  \item \textbf{ST463 (C6):} Read \texttt{/note="genotype: ST463"} from the MH734334.1 LOCUS metadata (submitter-provided; not independently re-typed because isolate WGS is not on SRA).
\end{itemize}

\subsection{Tools}
Python 3 stdlib only; NCBI BLAST+ (\texttt{blastn}, \texttt{makeblastdb}); \texttt{grep}/\texttt{awk} for parsing. Sequence content is deterministic by accession, so version drift is impossible for the numeric claims.

\section{Results vs paper}
\begin{longtable}{@{}p{0.04\textwidth}p{0.20\textwidth}p{0.15\textwidth}p{0.25\textwidth}p{0.08\textwidth}p{0.20\textwidth}@{}}
\toprule
\# & Paper claim & Paper value & Replication value & Match? & Note \\ \midrule
C1 & Plasmid length & 62{,}793~bp & \textbf{62{,}793~bp} & \checkmark exact & Directly from MH734334.1 \\
C2 & GC content & 58.8\% & \textbf{58.78\%} & \checkmark ($\pm$0.02\%) & Rounded, matches \\
C3 & \textit{bla}\textsubscript{KPC-2} present & Yes (PCR) & 293-aa ORF, \textbf{100.00\% (293/293)} identity to canonical KPC-2 & \checkmark exact & AZZ88873.1 confirms \\
C4 & Genetic env.\ near \textit{bla}\textsubscript{KPC-2} & Tn3 + ISKpn upstream; ISKpn6/IS26 downstream & Tn3 resolvase 15{,}740--16{,}297~(+); hypothetical CDS 16{,}420--17{,}400 (ISKpn8 position); \textit{bla}\textsubscript{KPC-2} 17{,}676--18{,}557~(+); downstream hypothetical CDSs 18{,}807--20{,}289 (ISKpn6/IS26 positions); KlcA 20{,}417--20{,}842 & \checkmark structurally & Some IS ORFs annotated as ``hypothetical protein'' in submitter GenBank; positions match paper schema \\
C5 & Novel backbone & ``novel'' & Shares \textbf{51{,}587 / 62{,}793 = 82.15\%} of length with p14057 at \textbf{98.70\%} weighted identity; $\sim$11.2~kb ($\sim$18\%) unique to pPA1011 & \textcolor{orange}{$\triangle$} partial & Novel \emph{configuration/insertion}, not unrelated plasmid family \\
C6 & PA1011 is ST463 & ST463 & ST463 (per MH734334.1 record note) & \checkmark provenance & Would need WGS + MLST scheme to independently verify \\
\bottomrule
\end{longtable}

\section{Verdict}
\textbf{REPLICATED} for the hard, numeric claims (C1, C2, C3, C4). Plasmid length, GC, presence of \textit{bla}\textsubscript{KPC-2} (100\% protein identity to canonical), and the local genetic environment near \textit{bla}\textsubscript{KPC-2} all reproduce exactly from the deposited GenBank record.

\textbf{PARTIAL / soft} on novelty (C5). The genetic-environment novelty ($\Delta$IS6--Tn3--ISKpn8--\textit{bla}\textsubscript{KPC-2}--ISKpn6--IS26) is defensible for \textit{P.~aeruginosa}. The backbone-novelty claim is only weakly supported: pPA1011 shares $\sim$82\% of its sequence at $\sim$98.7\% identity with p14057 (KY296095), so it is more accurately a \emph{variant} of a p14057-family plasmid with $\sim$11~kb of additional/rearranged content, rather than an unrelated new plasmid family.

\textbf{PROVENANCE-only} on C6 (ST463): the paper's MLST call is retained from the submitter's GenBank note.

\section{GENUINE CRITIQUE}
This section is a deliberately adversarial reading of both the target paper and this replication, so that reviewers can see what was \emph{not} strengthened.

\subsection{Limitations of the target paper}
\begin{enumerate}[leftmargin=1.4em]
  \item \textbf{``Novel'' is oversold.} The paper describes pPA1011 as a ``novel plasmid'', but our BLAST-based comparison against a single prior \textit{P.~aeruginosa} KPC-2 plasmid (p14057, KY296095) already recovers $\sim$82\% of pPA1011 at 98.7\% identity. The novelty story would be better phrased as ``novel IS-mediated \textit{bla}\textsubscript{KPC-2} context grafted onto a p14057-family backbone'', which is a materially weaker claim than a truly unrelated plasmid.
  \item \textbf{No raw-read deposit.} The Illumina NextSeq~500 and PacBio RSII reads underlying the assembly are not deposited under an accessible SRA/ENA accession for this study. This forecloses independent reassembly, error-model auditing, and coverage-based confirmation of the IS boundaries.
  \item \textbf{No isolate WGS.} Only the plasmid is deposited; the chromosomal genome of PA1011 is not on SRA under this study. Consequently, the ST463 MLST call cannot be independently re-typed, and one cannot check for chromosomal copies of \textit{bla}\textsubscript{KPC-2} or resistance-context context (e.g., a chromosomal Tn3 island).
  \item \textbf{IS-family annotation via ``hypothetical protein''.} The submitter's GenBank annotation labels several transposase/IS ORFs generically as ``hypothetical protein'' rather than by IS family (ISKpn8, ISKpn6, IS26). The paper's schema is only \emph{consistent with}, not \emph{proven by}, the deposited record; a definitive assignment would require ISfinder / ISEScan re-annotation.
  \item \textbf{No functional resistance data at the sequence level.} PCR + sequence presence of \textit{bla}\textsubscript{KPC-2} does not by itself confirm functional expression at the reported MIC levels; the paper's phenotype-to-genotype linkage relies on the clinical microbiology of the isolate, which is not sequence-checkable.
\end{enumerate}

\subsection{Limitations of this replication}
\begin{enumerate}[leftmargin=1.4em]
  \item \textbf{No independent assembly.} We accept MH734334.1 as ground truth. Any assembly artifact (mis-joined repeats, collapsed IS copies, incorrect circularisation) would silently propagate through C1--C4.
  \item \textbf{Single comparator for C5.} Novelty is assessed against \emph{one} closely related plasmid (p14057). A proper backbone-novelty test would BLAST against all NCBI \textit{P.~aeruginosa} KPC-carrying plasmids (or nucleotide-nr with taxonomic filter) and report the maximum matching backbone. This replication does not do that.
  \item \textbf{C4 is coordinate-consistency, not IS-typing.} We show that CDS/repeat features at the reported positions match the paper's schematic, but we do not run ISfinder/ISEScan to confirm the IS-family identities of the ``hypothetical protein'' ORFs.
  \item \textbf{No re-typed MLST for C6.} Same reason as the paper: no accessible isolate WGS.
  \item \textbf{No PCR replication for C3.} We upgrade C3 to a sequence-level protein-identity check, which is stronger than PCR, but we do not verify that the deposited sequence corresponds to the same physical isolate as the paper's PCR (this requires trust in the deposit chain).
\end{enumerate}

\subsection{What would make the replication stronger}
Reassemble from raw reads (requires SRA deposit that does not exist); run ISfinder/ISEScan on MH734334.1 to definitively type the IS ORFs; BLAST pPA1011 against all NCBI \textit{P.~aeruginosa} KPC-2 plasmids to properly bound the backbone-novelty claim; obtain the isolate WGS to independently re-type ST463 and check for chromosomal \textit{bla}\textsubscript{KPC-2}.

\section{Caveats}
This replication does \textbf{not} reassemble raw reads, does \textbf{not} independently re-type PA1011 by MLST, and does \textbf{not} re-run PCR (C3 is instead confirmed at the sequence level, which is a strictly stronger check). It does \textbf{not} claim the pPA1011 backbone is unique; on the contrary, it tightens ``novel'' into ``variant of a p14057-family plasmid with $\sim$11~kb of distinguishing content and a novel IS-mediated \textit{bla}\textsubscript{KPC-2} context''.

\end{document}
